\subsection{分解}

在下文中，总是将模与一个复形等同起来，该复形以该模为零次分量，且所有其他分量为零。

定义1. --- 设 M 为一个 A-模。M 的一个左分解是一个偶对 (P, p)，其中 P 是一个在右侧为零的复形，且 p : P → M 是一个拟同构。M 的一个右分解是一个偶对 (e, E)，其中 E 是一个在左侧为零的复形，且 e : M → E 是一个拟同构。

称分解(P, p)（相应地(e, E)）的长度为复形P（相应地E）的长度。若(P, p)和(P', p')（相应地(e, E)和(e', E')）是M的两个左（相应地右）分解，则一个复形态射 \(f: P \to P'\) 使得 \(p' \circ f = p\) (相应地， \(g: E \to E'\) 使得 \(g \circ e = e'\) ）被称为分解的态射。

命题2. --- 设 P 为一个在右侧为零的复形，且设 p : P → M 为一个态射。要使 (P, p) 成为 M 的一个左分解，必要且充分条件是序列

\[
\dots \longrightarrow \mathrm{P} _ {n} \xrightarrow {d _ {\mathrm{P}}} \mathrm{P} _ {n - 1} \longrightarrow \dots \quad \longrightarrow \mathrm{P} _ {1} \xrightarrow {d _ {\mathrm{P}}} \mathrm{P} _ {0} \xrightarrow {p _ {0}} \mathrm{M} \longrightarrow 0\tag{1}
\]

是正合的。

事实上，说 \(p: \mathbf{P} \to \mathbf{M}\) 是拟同构，意味着 \(\mathrm{H}_{i}(\mathbf{P}) = 0\) （对 \(i > 0\) ），且 \(p_{0}\) 诱导了从 Coker ( \(d_{\mathbf{P}}: \mathbf{P}_{1} \to \mathbf{P}_{0}\) ) 到 \(\mathbf{M}\) 上的一个同构 .

的同构。类似地：

命题2bis. --- 设 E 是一个在左侧为零的复形，并设 e : M → E 是一个态射。要使 (e, E) 成为 M 的一个右分解，必要且充分条件是序列

\[
0 \longrightarrow \mathrm{M} \stackrel {e _ {0}} {\longrightarrow} \mathrm{E} ^ {0} \stackrel {d _ {\mathrm{E}}} {\longrightarrow} \mathrm{E} ^ {1} \longrightarrow \dots \longrightarrow \mathrm{E} ^ {n} \stackrel {d _ {\mathrm{E}}} {\longrightarrow} \mathrm{E} ^ {n + 1} \longrightarrow \dots\tag{1 bis}
\]

是正合的。

为方便起见，人们常说序列 (1)（相应地 (1 bis)）是 M 的左（相应地右）分解。

定义2. --- A-模 M 的一个投射（相应地自由、相应地平坦）分解是 M 的一个左分解 (P, p)，使得复形 P 是投射的（相应地自由的、相应地平坦的）（X，第 25 页）。M 的一个内射分解是 M 的一个右分解 (e, E)，使得复形 E 是内射的（同上）。

例子. --- 1) 假设环 A 是主理想整环；设 M 是一个 A-模，并设 \((x_{i})_{i\in I}\) 是 M 的一个生成族。设 \(L_{0}\) 是自由模 \(\mathbf{A}^{(I)}\) , \((e_{i})\) 为其典范基，并定义 \(p:L_{0}\to\mathbf{M}\) 为 \(p(e_{i})=x_{i}\) 。态射 \(p\) 是满射的，且其核 \(L_{1}\) 由 VII, § 3, 定理 1 的推论 2 知是一个自由 A-模，因此正合列

\[
0 \to \mathrm{L} _ {1} \to \mathrm{L} _ {0} \xrightarrow {p} \mathrm{M} \to 0
\]

是 M 的长度为 1 的自由分解。若 I 是有限的，则 \(L_{0}\) 与 \(L_{1}\) 是有限生成的。

\begin{enumerate}
\def\labelenumi{\arabic{enumi})}
\setcounter{enumi}{1}
\tightlist
\item
  假设 A 是交换的；设 E 是一个 A-模，u 是 E 的一个自同态。记 \(E_{u}\) 为通过赋予 E 由下式定义的结构而得到的 A{[}X{]}-模
\end{enumerate}

\[
(p, x) \mapsto p (u) (x) \quad \text { for } \quad p \in \mathrm{A} [ \mathrm{X} ] \quad \text { and } \quad x \in \mathrm{E}.
\]

根据III，第106页，我们有一个正合列：

\[
0 \to \mathrm{A} [ \mathrm{X} ] \otimes_ {\mathrm{A}} \mathrm{E} \stackrel {\psi} {\to} \mathrm{A} [ \mathrm{X} ] \otimes_ {\mathrm{A}} \mathrm{E} \stackrel {\varphi} {\to} \mathrm{E} _ {u} \to 0
\]

其中 \(\varphi(p \otimes x) = p.x\) 与 \(\psi(p \otimes x) = Xp \otimes x - p \otimes u(x)\) （当 \(p \in A[X]\) 与 \(x \in E\) ）。这个正合列是 \(E_u\) 的长度为 1 的自由（相应地，投射，相应地，有限生成）分解，如果 E 是自由（相应地，投射，相应地，有限生成）A-模。3) 如果 A 是主理想整环，则正合列

\[
0 \rightarrow \mathrm{A} \rightarrow \mathrm{K} \rightarrow \mathrm{K/A} \rightarrow 0
\]

是 A-模 \(\mathbf{A}_s\) 的长度为 1 的内射分解（X，第 18 页，例 1）。

命题 3. --- 设 \(f: \mathbf{M}' \to \mathbf{M}\) 是 A-模的一个同态， \(p': \mathbf{P}' \to \mathbf{M}'\) 是到 \(\mathbf{M}'\) 的、来自一个在右侧为零且投射的复形 \(\mathbf{P}'\) 的一个态射，且 \(p: \mathbf{P} \to \mathbf{M}\) 是 M 的一个左分解。存在一个复形态射 \(\tilde{f}: \mathbf{P}' \to \mathbf{P}\) ，且在同伦意义下唯一，使得 \(p \circ \tilde{f} = f \circ p'\) .

考虑如下定义的复形 \(\overline{\mathbf{P}}\) ： \(\overline{\mathbf{P}}_n = \mathbf{P}_n\) 当 \(n \neq -1\) , \(\overline{\mathbf{P}}_{-1} = \mathbf{M}\) , \(d_{\overline{\mathbf{P}},n} = d_{\mathbf{P},n}\) 当 \(n \neq 0, -1\) , \(d_{\overline{\mathbf{P}},0} = p_0\) , \(d_{\overline{\mathbf{P}},-1} = 0\) ，以及类似定义的复形 \(\overline{\mathbf{P}}'\) 。将命题 1 a) 应用于复形 \(\overline{\mathbf{P}}\) 与 \(\overline{\mathbf{P}}'\) ，取 \(r=0\) , \(u_i = 0\) （当 \(i < -1\) ）与 \(u_{-1} = f\) ，我们得到命题 3。

推论。--- 设 (P, p) 与 (P', p') 是 M 的两个投射分解。存在一个在同伦意义下唯一的同伦等价 \(\alpha: P' \to P\) ，使得 \(p \circ \alpha = p'\) .

事实上，存在一个态射 \(\alpha: P' \to P\) (相应地， \(\beta: P \to P'\) ）使得 \(p \circ \alpha = p'\) (相应地， \(p' \circ \beta = p\) ）。由于 \(p \circ \alpha \circ \beta = p\) (相应地， \(p' \circ \beta \circ \alpha = p'\) ), \(\alpha \circ \beta\) 同伦于 \(1_P\) (相应地， \(\beta \circ \alpha\) 同伦于 \(1_{P'}\) ).

命题 3 bis. --- 设 \(g: \mathbf{N} \to \mathbf{N}'\) 是 A-模的一个同态， \(e': \mathbf{N}' \to \mathbf{E}'\) 是 \(\mathbf{N}'\) 到一个在左侧为零且内射的复形 \(\mathbf{E}'\) 中的态射，且 \(e: \mathbf{N} \to \mathbf{E}\) 是 \(\mathbf{N}\) 的一个右分解。存在一个复形态射 \(\tilde{g}: \mathbf{E} \to \mathbf{E}'\) ，且在同伦意义下唯一，使得 \(\tilde{g} \circ e = e' \circ g\) .

此证明类似于命题3，借助命题1 b) 完成。

推论。--- 设 \((e, E)\) 与 \((e', E')\) 是 N 的两个内射分解；存在一个同伦等价 \(\alpha : E \to E'\) ，且在同伦意义下唯一，使得 \(\alpha \circ e = e'\) .

